\documentclass[a4paper,10pt,twoside]{report}
\usepackage{geometry}\geometry{a4paper,top=3.5cm,bottom=3.5cm,%
left=2.5cm,right=2.5cm,heightrounded,bindingoffset=0mm}
\usepackage[T1]{fontenc}
\usepackage[utf8]{inputenc}
\usepackage[italian]{babel}
\usepackage{graphicx}
\usepackage[export]{adjustbox}%Per il Frame attrono le immagini e il valign
\usepackage{subfig}
\usepackage{amsmath,amsfonts,amssymb,braket,mathrsfs}
\usepackage{float}
\usepackage{tabularx,booktabs}
\usepackage{hyperref}
\usepackage{epsfig}
\usepackage{pdfpages} %Per gli allegati
%\usepackage{minipage}
\usepackage[output-decimal-marker={,}]{siunitx}
\DeclareSIUnit \days {gg}
%ILe tre righe sotto ervono per mettere il grassetto dentro la tabelle siunitex usando \B, il rosso usando \RED, e il verde usando \GREEN
\sisetup{detect-weight,mode=text}
\usepackage{etoolbox}
\newrobustcmd\B{\DeclareFontSeriesDefault[rm]{bf}{b}\bfseries}
\newrobustcmd\RED{\DeclareFontSeriesDefault[rm]{bf}{b}\color{red}}
\newrobustcmd\GREEN{\DeclareFontSeriesDefault[rm]{bf}{b}\color{myGreen}}
%%
\usepackage{tikz}
\usepackage{pgfplots,pgfplotstable}
%\pgfplotsset{compat=1.15} %indica la versione da utilizzare per pgfplot
\usetikzlibrary{patterns} % per il tratteggio
\usepgfplotslibrary{groupplots}
\pgfplotsset{compat=newest}
%\usepackage{stanli}
\usepackage{xspace}% per lo spazio intelligente
\newcommand{\e}{\`E\xspace}  %E'
\usepackage{titlesec} % per formato custom dei titoli dei capitoli
%\usepackage{sideways}%%%
% redefinizione del formato del titolo del capitolo
      % da formato
      %   Capitolo X
      %   Titolo capitolo
      %   a formato
      %      Titolo capitolo 
	\titleformat{\chapter}
        {\normalfont\Huge\bfseries}{}{0em}{}
	\titlespacing*{\chapter}{0pt}{0in}{0.02in}
	\titlespacing*{\section}{0pt}{0.2in}{0.02in}
	\titlespacing*{\subsection}{0pt}{0.10in}{0.02in}
%serve per la didascalia di tabelle e figure:
\usepackage{caption}
\captionsetup{tableposition=top,figureposition=bottom,font=small}\captionsetup{format=hang,labelfont={bf,color=pantone186}} %didascalie a più righe allineate e il nome in grassetto
%non viene allineato a sinistra se la didascalia è corta una sola riga. PERCHé??
\usepackage{xcolor}
%serve per mettere il codice con lo sfondo grigio chiaro
\definecolor{pantone186}{RGB}{206, 17, 38} %il colore del logo UNITN
\definecolor{myGray}{gray}{0.5} %più basso più scuro è
\definecolor{myGreen}{rgb}{0.0, 0.5, 0.0}
\usepackage{listings} 
\lstset{basicstyle=\small\ttfamily,
backgroundcolor=\color{black!10},%
boxpos=c,%
stringstyle=\itshape,		
lineskip=3pt,%
numbers=left,
numberstyle=\footnotesize,}
\usepackage{lscape}
\usepackage{multirow}
\usepackage{import}
\begin{document}
\pagestyle{plain}
\thispagestyle{empty}
\begin{center}
  \begin{figure}[H]
    \centerline{\psfig{file=IMG/logo_unitn_black_centerNEW.eps,
                        width=0.8\textwidth,trim = 0 0.9cm 0 0.5cm}}
  \end{figure}
  \textcolor{pantone186}{\noindent\rule{\textwidth}{.5pt}}

  \Large\textsc{Dipartimento di Ingegneria Civile, Ambientale e Meccanica\\}
  \Large{Corso di Laurea in Ingegneria Civile
  }

  \vspace{3.2 cm} 
  %\Large\textsc{Elaborato finale\\} 
  %\vspace{1 cm} 
  \Huge\textsc{Relazione di Tecnica delle Costruzioni\\}
  
  \vspace{0.2 cm}
  \Large{\it{Progetto e verifica di alcuni elementi strutturali in C.A.\\
  di un edificio multipiano con struttura a telaio}}


  \vspace{4 cm} 
  \begin{tabular*}{\textwidth}{ l @{\extracolsep{\fill}} r }
  \Large\textsc{Docenti} & \Large\textsc{Studente}\\
  \Large{Nadia Baldassino}& \Large{Nicola Meoli 225077}\\
  \Large{Stefano Gasperetti}& \\

  	
  	
  \end{tabular*}

  \vspace{3.1cm} 
  \textcolor{pantone186}{\noindent\rule{\textwidth}{1pt}}
    
  \Large{Anno accademico 2021/22}
  
\end{center}

\tableofcontents
%\setcounter{page}{1}
%Tabelle e figure sulla stessa pagina:
%Le aggiunge all'indice. phantomsection serve per non far casini con hyperref
\clearpage
\begingroup
   %\let\cleardoublepage\relax  % book
    \let\clearpage\relax        % report
        \listoftables
        \phantomsection
        \addcontentsline{toc}{chapter}{Elenco delle tabelle}
\endgroup
        %
        \clearpage
\begingroup
        \let\clearpage\relax
        \listoffigures
        \phantomsection
        \addcontentsline{toc}{chapter}{Elenco delle figure}
\endgroup
%
%%%%%%%%%%
%Comandi aggiunti:
%%%%%%%%%%
\newcommand{\red}[1]{\textcolor{pantone186}{#1}}
%

\chapter{Introduzione}
La presente relazione riguarda il progetto e la verifica strutturale di alcuni elementi componenti un edificio composto da 4 piani di cui 3 fuori terra sito in Provincia di Trento ad un'altitudine di \SI{300}{\meter} sopra il livello del mare.
Le dimensioni in pianta sono di circa $28 \times \SI{42}{\metre}$ per una altezza totale di \SI{12.45}{\meter} di cui \SI{9.70}{\meter} fuori terra. 
La cubatura totate è perciò di circa \SI{14650}{\metre\cubed}.

\paragraph{Composizione strutturale dell'edificio}
La struttura dell'edificio è a telaio in calcestruzzo armato, da progetto preliminare caratterizzata da pilastri quadrati di sezione \SI{30}{\centi\meter} e da travi sia a spessore che non, rispettivamente larghe \SI{60}{\centi\meter} ed alte come lo spessore del solaio, e $30 \times \SI{50}{\centi\metre}$.

I solai strutturali sono di tipo a lastre tralicciate Predalle tra il piano interrato e il piano terra e nel solaio di copertura. 
Il peso ultimato di tale solaio di \SI{3.6}{\kilo\newton\per\square\meter}.
I solai tra piano terra e piano primo e tra piano primo e piano secondo sono di tipo a travetti e latero-cemento.
Il peso ultimato di tale solaio di \SI{3.20}{\kilo\newton\per\square\meter}.

I solai dei piano intermedi sono costituiti da un pacchetto con calcestruzzo alleggerito, massetto di allettamento, intonaco e pavimento in ceramica.
Il solaio nella zona del terrazzo è costituito da isolante, impermealizzazione, massetto in calcestruzzo e pavimento.
Il solaio di copertura è costituito da isolante, massetto in calcestruzzo alleggerito, impermealizzazione e ghiaia.

Le pareti divisorie interne sono costituite da muratura in laterizio e intonaco in ambo i lati.
Le pareti perimetrali sono costituite da muratura in laterizio, cappotto esterno e intonaco.

Il piano terra è adibito a negozi, il piano primo ad uffici aperti al pubblico, il piano secondo a civile abitazione.
Il piano interrato è adibito a garage.

\paragraph{Oggetto della relazione}
Gli elementi che si andranno ad analizzare saranno una trave ed un solaio del primo piano e il pilastro P27. 
Tali elementi sono ben evidenziati nelle piante a seguire insieme alle sollecitazioni in essi agenti.

Subito dopo vengono elencati normative e materiali impiegati nell'analisi strutturale.

Seguono poi i capitoli di calcolo relativi agli elementi in esame.
La trave verrà analizzata per sollecitazioni agli SLU e SLE, solaio e pilastro per quelle agli SLU.

In allegato al presente documento sono riportati degli elaborati grafici eseguiti in software CAD.

\end{document} 